\section{Resumen Ejecutivo}\label{resumen-ejecutivo}

Este informe documenta la containerizacion de InvenTrack, su despliegue en un cluster Kubernetes local mediante Minikube y la auditoria de seguridad basada en PTES. La aplicacion se desplego y se valido bajo el dominio de laboratorio \texttt{conjunta3p.espe.edu.ec}, resuelto localmente hacia el entorno controlado.

El proceso se ejecuta con separacion de entornos: Docker, kubectl y Minikube se operan desde WSL Ubuntu; Kali Linux WSL se reserva para las fases de pentesting y recoleccion de evidencias tecnicas.

\section{Objetivos}\label{objetivos}

\begin{itemize}
\tightlist
\item
  Crear imagenes Docker seguras para backend, frontend y MySQL.
\item
  Crear manifiestos Kubernetes para namespace, deployments, services, ConfigMap, Secret, PVC e Ingress.
\item
  Desplegar InvenTrack en Minikube y publicar la aplicacion mediante Ingress.
\item
  Resolver el dominio \texttt{conjunta3p.espe.edu.ec} hacia el laboratorio local.
\item
  Ejecutar y documentar las siete fases PTES.
\item
  Clasificar hallazgos por severidad y proponer mitigaciones accionables.
\end{itemize}

\section{Alcance}\label{alcance}

El alcance autorizado se limita a la instancia local de InvenTrack desplegada para esta practica.

Activos incluidos:

{\def\LTcaptype{none} % do not increment counter
\captionof{table}{Activos incluidos en el alcance}\label{tab:activos-alcance}
\begin{longtable}[]{@{}ll@{}}
\toprule\noalign{}
Activo & Alcance \\
\midrule\noalign{}
\endhead
\bottomrule\noalign{}
\endlastfoot
Frontend React/Vite & Incluido \\
Backend Node.js/Express & Incluido \\
API REST\texttt{/api} & Incluido \\
Cargas\texttt{/uploads} & Incluido \\
Base de datos MySQL del laboratorio & Incluida indirectamente desde la aplicacion \\
Autenticacion JWT y roles & Incluidos \\
Chatbot IA integrado & Incluido solo como endpoint de InvenTrack \\
\end{longtable}
}

Exclusiones:

\begin{itemize}
\tightlist
\item
  Infraestructura real de \texttt{espe.edu.ec}.
\item
  Otros subdominios institucionales.
\item
  Ataques de denegacion de servicio.
\item
  Extraccion completa de datos.
\item
  Pruebas directas contra la API de Groq.
\item
  Publicacion de secretos reales, tokens o hashes completos.
\end{itemize}

\section{Reglas de Juego}\label{reglas-de-juego}

\begin{itemize}
\tightlist
\item
  Todas las pruebas activas se ejecutan con bajo volumen y datos de laboratorio.
\item
  Antes de escanear desde Kali, se verificara que \texttt{conjunta3p.espe.edu.ec} resuelve al laboratorio local.
\item
  Los Secrets reales se crean en Kubernetes y no se versionan.
\item
  Toda evidencia con credenciales, tokens, cookies, API keys o hashes completos sera redactada antes de incluirse en el PDF.
\item
  Las pruebas se detendran si el dominio resuelve a una IP publica o no esperada.
\end{itemize}

\section{Arquitectura Objetivo}\label{arquitectura-objetivo}

\begin{tcolorbox}[codebox]
\begin{Verbatim}[fontsize=\scriptsize,breaklines=true,breakanywhere=true]
Navegador / Kali
      |
      v
conjunta3p.espe.edu.ec
      |
      v
Ingress Nginx - Minikube
      |
      +--> /                 frontend-service
      +--> /api, /uploads    backend-service
                                |
                                v
                            mysql-service
                                |
                                v
                            MySQL + PVC
\end{Verbatim}
\end{tcolorbox}

\section{Base teorica y adaptacion Kubernetes}\label{base-teorica-y-adaptacion-kubernetes}

La base autoritativa de la practica es implement.md; el plan operativo es PLAN\_IMPLEMENTACION\_PTES\_KALI\_WSL.md. Como apoyo teorico se incorporo guia-pentest-k8s.md, especialmente para analizar Ingress, namespaces, Services ClusterIP, NetworkPolicies, Secrets, PVC, securityContext, ServiceAccounts y fronteras de confianza.

La guia teorica tambien contiene tecnicas para API Server, kubelet, etcd, Dashboard, tokens de ServiceAccount, extraccion de Secrets, RCE, escape de contenedor y movimiento lateral. Esas tecnicas no se ejecutan en esta practica: implement.md y la autorizacion P01-01 limitan la auditoria a la aplicacion InvenTrack publicada bajo conjunta3p.espe.edu.ec. Tambien quedan fuera MySQL directo, red interna del cluster, DoS y cualquier infraestructura real de ESPE.

La guia se usa, por tanto, para fundamentar el modelo de amenazas y los controles Kubernetes, no para ampliar el alcance autorizado.
